\documentclass[a4paper,12pt]{article}

\usepackage{geometry}
\geometry{top=20mm,bottom=20mm,left=25mm,right=25mm}

\usepackage[utf8]{inputenc}
\usepackage[T2A]{fontenc}
\usepackage[russian]{babel}

\usepackage{xcolor}
\usepackage{listings}
\usepackage{hyperref}

\definecolor{codebg}{RGB}{248,248,248}
\definecolor{codeframe}{RGB}{210,210,210}
\definecolor{codekeyword}{RGB}{0,76,153}
\definecolor{codecomment}{RGB}{0,128,0}
\definecolor{codestring}{RGB}{163,21,21}
\definecolor{codenumber}{RGB}{120,120,120}

\lstdefinestyle{common}{
    backgroundcolor=\color{codebg},
    frame=single,
    rulecolor=\color{codeframe},
    basicstyle=\ttfamily\small,
    keywordstyle=\color{codekeyword}\bfseries,
    commentstyle=\color{codecomment}\itshape,
    stringstyle=\color{codestring},
    numberstyle=\tiny\color{codenumber},
    numbers=left,
    numbersep=10pt,
    breaklines=true,
    showstringspaces=false,
    keepspaces=true,
    columns=fullflexible,
    tabsize=4
}

\lstdefinestyle{cppstyle}{
    style=common,
    language=C++
}

\lstdefinestyle{pythonstyle}{
    style=common,
    language=Python
}




\begin{document}

\begin{titlepage}
    \centering
    {\large МИНОБРНАУКИ РОССИИ\par}
    \vspace{0.5em}
    {\large САНКТ-ПЕТЕРБУРГСКИЙ ГОСУДАРСТВЕННЫЙ\par}
    {\large ЭЛЕКТРОТЕХНИЧЕСКИЙ УНИВЕРСИТЕТ\par}
    {\large «ЛЭТИ» ИМ. В.И. УЛЬЯНОВА (ЛЕНИНА)\par}
    \vspace{1em}
    {\large Кафедра Информационных систем\par}

    \vspace{6em}

    {\LARGE \textbf{Практическая работа}\par}
    \vspace{0.8em}
    {\large по дисциплине «Теория принятия решений»\par}
    \vspace{0.8em}
    {\large Тема: "Применение методов динамического программирования для решения задачи поиска оптимального плана заказа станков"\par}

    \vfill

    \begin{flushleft}
        Студент гр. 3376 \hfill Михайлов Н.Ф.\\[0.5em]
        Преподаватель   \hfill Пономарёв А.В.
    \end{flushleft}

    \vspace{3em}

    Санкт-Петербург\par
    2026
\end{titlepage}

\tableofcontents
\newpage

\section{Введение}
Целью в данной практической работе является освоение метода динамического программирования
в рамках задачи о поиске оптимального плана заказа станков. \\
Для этого нам необходимо
\begin{enumerate}
    \item Формализовать задачу
    \item Составить уравнение Беллмана
    \item Реализовать алгоритм (прямой или обратной прогонки)
    \item Протестировать алгоритм и проанализировать его итоговые данные
\end{enumerate}

\section{Постановка задачи}
Склад пункта реализации станков имеет вместимость 20 единиц. Пополнение склада возможно только первого числа каждого месяца.
Станки привозят автотранспортом (1 рейс), причем стоимость рейса складывается из постоянных затрат 50 денежных единиц (ДЕ)
и затрат на доставку каждого станка 10 ДЕ. За один рейс (и, следовательно, в месяц) может быть доставлено не более 5 станков.
\\
Затраты на хранение станка в течение месяца составляют 10 ДЕ.
\\
Ожидаемый спрос на станки приведен в табл. 1. Отсутствие требуемого количества станков на складе недопустимо.
В начальный момент на складе находится 1 станков.
\begin{enumerate}
	\item Определить план заказов, минимизирующий стоимость
	\item Определить границы изменения стоимости хранения станка, в которых найденный план является оптимальным
	\item Определить границы изменения постоянных затрат на совершение рейса, в которых най-
денный план является оптимальным
	\item Определить план заказов, минимизирующий стоимость, при условии, что к концу периода
на складе должно остаться 2 станков
\end{enumerate}
\section{Формализация задачи и вывод уравнения Беллмана}

\subsection{Этапы, состояние и управление}
В качестве этапов рассматриваются месяцы планового периода, то есть
\[
i = 1,\dots,6.
\]

Состояние системы на шаге $i$ определяется количеством станков на складе в начале месяца. Обозначим это состояние через
\[
S_i.
\]

Управлением на шаге $i$ является количество станков, заказываемых в начале месяца:
\[
u_i.
\]

Так как за один месяц может быть выполнен только один рейс, а за один рейс можно доставить не более 5 станков, управление ограничено:
\[
0 \le u_i \le 5.
\]

Переход системы из состояния $S_i$ в состояние $S_{i+1}$ задаётся уравнением баланса:
\[
S_{i+1} = S_i + u_i - d_i,
\]
где $d_i$ --- спрос на станки в месяце $i$.

Таким образом, состояние на следующем шаге зависит от текущего запаса на складе, принятого решения о заказе и спроса в рассматриваемом месяце.

\subsection{Ограничения}
При решении задачи необходимо учитывать следующие ограничения.

\begin{enumerate}
    \item Вместимость склада не должна превышаться:
    \[
    0 \le S_i \le 20.
    \]

    \item Дефицит станков недопустим, поэтому имеющегося запаса и объёма заказа должно хватать на покрытие спроса:
    \[
    S_i + u_i \ge d_i.
    \]

    \item За один месяц можно заказать не более 5 станков:
    \[
    0 \le u_i \le 5.
    \]

    \item Для подзадачи 4 добавляется дополнительное ограничение на конечный остаток:
    \[
    S_7 = 2.
    \]
\end{enumerate}

Эти ограничения определяют множество допустимых решений, среди которых выбирается решение с минимальной стоимостью.
\subsection{Уравнение Беллмана}
Пусть
\[
W_i(S_i)
\]
--- минимальные суммарные затраты от месяца $i$ до конца периода, если в начале месяца $i$ на складе находится $S_i$ станков.

Тогда уравнение Беллмана для данной задачи имеет вид:
\[
W_i(S_i)=
\min_{u_i \in U_i(S_i)}
\left[
c(u_i) + h S_{i+1} + W_{i+1}(S_{i+1})
\right],
\]
где:
\begin{itemize}
    \item $h$ --- стоимость хранения одного станка;
    \item $S_{i+1} = S_i + u_i - d_i$;
    \item $U_i(S_i)$ --- множество допустимых управлений;
    \item $c(u_i)$ --- стоимость поставки в месяце $i$.
\end{itemize}

Стоимость поставки определяется следующим образом:
\[
c(u_i)=
\begin{cases}
0, & u_i = 0, \\
50 + 10u_i, & u_i > 0.
\end{cases}
\]

Тогда уравнение Беллмана можно записать более подробно:
\[
W_i(S_i)=
\min_{u_i}
\left[
\begin{cases}
0, & u_i=0, \\
50 + 10u_i, & u_i>0
\end{cases}
+ 10S_{i+1} + W_{i+1}(S_{i+1})
\right].
\]

Граничное условие для базовой задачи:
\[
W_7(S_7)=0.
\]

Для подзадачи 4, где требуется получить конечный остаток 2 станка:
\[
W_7(S_7)=
\begin{cases}
0, & S_7=2, \\
+\infty, & S_7 \ne 2.
\end{cases}
\]

Это уравнение лежит в основе реализованного рекурсивного алгоритма.

\section{Решение подзадач}
\subsection{Подзадача 1}
В первой подзадаче требуется определить оптимальный план заказов, минимизирующий общую стоимость поставки и хранения.

Программа последовательно рассматривает все допустимые объёмы заказа в каждом месяце, вычисляет стоимость текущего решения и рекурсивно добавляет минимальную стоимость на оставшемся горизонте планирования. После нахождения минимальной стоимости восстанавливается сам план заказов по месяцам.

В результате был получен следующий оптимальный план:
\[
(2, 5, 5, 5, 0, 0).
\]

Минимальная суммарная стоимость:
\[
450.
\]

Это решение означает, что выгодно сделать поставки в первые четыре месяца, а затем использовать накопленный запас на складе.
\subsection{Подзадача 2}
Во второй подзадаче требуется определить границы изменения стоимости хранения одного станка, при которых найденный в первой подзадаче план остаётся оптимальным.

Для этого фиксируется базовый оптимальный план, после чего параметр стоимости хранения последовательно уменьшается и увеличивается. Для каждого нового значения параметра задача решается заново, а полученный план сравнивается с базовым.

В результате установлено, что базовый план остаётся оптимальным при стоимости хранения в диапазоне:
\[
[0; 12].
\]

Следовательно, если стоимость хранения не превышает 12 денежных единиц за станок в месяц, структура оптимального плана не изменяется.
\subsection{Подзадача 3}
В третьей подзадаче исследуется влияние постоянных затрат на совершение рейса на оптимальность базового плана.

Аналогично подзадаче 2, фиксируется найденный план и выполняется параметрическое изменение постоянной составляющей стоимости рейса. При каждом новом значении параметра задача решается заново и сравнивается с исходным решением.

Нижняя граница изменения постоянных затрат составляет:
\[
41.
\]

Верхняя граница отсутствует. Это объясняется тем, что найденный базовый план использует минимально возможное число рейсов. Так как при увеличении постоянной стоимости рейса все планы с большим количеством рейсов становятся только менее выгодными, базовый план остаётся оптимальным и далее.

Таким образом, диапазон имеет вид:
\[
[41; +\infty).
\]
\subsection{Подзадача 4}
В четвёртой подзадаче требуется определить оптимальный план заказов при дополнительном условии:
\[
S_7 = 2.
\]

Для этого в граничное условие рекурсивного алгоритма вводится требование, чтобы к концу периода на складе осталось ровно 2 станка. Все остальные конечные состояния считаются недопустимыми.

В результате получен следующий план:
\[
(2, 5, 5, 3, 0, 4).
\]

Минимальная стоимость при этом составляет:
\[
500.
\]

По сравнению с базовой задачей стоимость возрастает, поскольку необходимо обеспечить дополнительный запас к концу периода.

\section{Сравнение скорость работы разных реализаций}
Для оценки эффективности были реализованы несколько версий алгоритма: на C++ и Python. Дополнительно была предусмотрена возможность запуска с мемоизацией и без неё.

Сравнение выполнялось с помощью сценария автоматического тестирования, который запускал каждую реализацию в одинаковых условиях и измерял время выполнения.

\subsection{Результаты измерений}
Для C++ и Python были получены следующие результаты:
\begin{itemize}
    \item C++ с мемоизацией --- $0.004$ с;
    \item C++ без мемоизации --- $0.007$ с;
    \item Python с мемоизацией --- $0.017$ с;
    \item Python без мемоизации --- $0.060$ с.
\end{itemize}

Из этих данных можно сделать следующие выводы:
\begin{itemize}
    \item для реализации на C++ использование мемоизации уменьшило время работы с $0.007$ с до $0.004$ с, то есть примерно на $42.9\%$;
    \item для реализации на Python использование мемоизации уменьшило время работы с $0.060$ с до $0.017$ с, то есть примерно на $71.7\%$;
    \item при использовании мемоизации Python работает примерно в $4.25$ раза медленнее C++, что соответствует увеличению времени примерно на $325\%$ относительно C++;
    \item без мемоизации Python работает примерно в $8.57$ раза медленнее C++, что соответствует увеличению времени примерно на $757\%$ относительно C++.
\end{itemize}

Следует отметить, что сама по себе рассматриваемая задача не является вычислительно очень сложной, поэтому все измеренные времена остаются малыми и на практике оба варианта работают быстро. Однако в относительном выражении разница оказывается весьма существенной.

Таким образом, мемоизация существенно повышает эффективность рекурсивного алгоритма динамического программирования, а выбор языка реализации также заметно влияет на производительность программы.
\section{Заключение}
В ходе работы была формализована задача оптимального планирования заказов станков и построено её решение методом динамического программирования на основе уравнения Беллмана.

Был реализован алгоритм, позволяющий:
\begin{itemize}
    \item находить оптимальный план заказов;
    \item исследовать устойчивость найденного решения к изменению параметров;
    \item учитывать дополнительное условие на конечный остаток на складе.
\end{itemize}

Проведённые вычисления показали, что метод динамического программирования позволяет эффективно решать задачу данного типа, а использование мемоизации заметно сокращает время работы алгоритма. Полученные результаты подтверждают применимость данного подхода для задач оптимального управления запасами.
\appendix
\pagebreak
\section{Листинг кода}

\subsection{Реализация на C++}
\lstinputlisting[style=cppstyle]{./lab1_listing.cpp}

\pagebreak
\subsection{Реализация на Python}
\lstinputlisting[style=pythonstyle]{./lab1_listing.py}


\end{document}
